% !TEX program = xelatex
\documentclass[UTF8,a4paper,11pt]{ctexart}

\usepackage[a4paper,margin=2.05cm,headheight=14pt]{geometry}
\usepackage{amsmath,amssymb,mathtools,bm}
\usepackage{booktabs,longtable,array,tabularx,multirow}
\usepackage{graphicx}
\usepackage{xcolor}
\usepackage{xurl}
\usepackage{hyperref}
\usepackage{pgfplots}
\usepackage{seqsplit}
\usepgfplotslibrary{groupplots}
\pgfplotsset{compat=1.18}

\definecolor{pyblue}{RGB}{36,91,163}
\definecolor{qudared}{RGB}{190,65,55}
\definecolor{stagegray}{RGB}{90,105,115}
\hypersetup{
  colorlinks=true,
  linkcolor=pyblue,
  urlcolor=pyblue,
  citecolor=pyblue,
  pdftitle={PyQCU/QUDA MultiGrid 分层性能与残差对照}
}
\urlstyle{tt}
\Urlmuskip=0mu plus 3mu
\newcommand{\pathref}[1]{\begingroup\Urlmuskip=0mu plus 3mu\path{#1}\endgroup}
\newcommand{\code}[1]{%
  \ifmmode
    \mathtt{\detokenize{#1}}%
  \else
    \begingroup\ttfamily
    \expandafter\seqsplit\expandafter{\detokenize{#1}}%
    \endgroup
  \fi}
\newcolumntype{L}[1]{>{\raggedright\arraybackslash}p{#1}}
\newcolumntype{Y}{>{\raggedright\arraybackslash}X}
\newcommand{\norm}[1]{\left\lVert#1\right\rVert}
\newcommand{\ip}[2]{\left\langle#1,#2\right\rangle}
\newcommand{\Yhat}{\widehat Y}
\setlength{\parindent}{2em}
\setlength{\parskip}{0.35em}
\emergencystretch=3em

\title{PyQCU 与 QUDA MultiGrid：可开关分层性能测试、\texorpdfstring{$r$}{r} 残差记录与三层加速}
\author{PyQCU 工程分析记录}
\date{2026 年 9 月 16 日}

\begin{document}
\maketitle

\begin{abstract}
本文在 PyQCU 与 QUDA 1.1.0 的 MultiGrid 求解路径中加入默认关闭、可显式开启的
分层性能与残差 trace，并在统一 gauge、右端项、canonical null vectors、精度、
block、odd--odd MATPC parity、停机容差和计时边界后完成对照。埋点不扩展公共
ABI：PyQCU 使用 \code{PYQCU_STRICT_TRACE_FILE}，QUDA 使用
\code{QUDA_MG_TRACE_FILE}。正式大格 $16\times32\times32\times48$ 的 c64
三层测试中，PyQCU steady solve 中位数为 $0.653995$ s，QUDA 为 $2.865627$ s，
PyQCU/QUDA 加速比为 $4.3817$；五次样本的最终全算符相对真残差分别为
$4.90135\times10^{-7}$ 与 $7.36873\times10^{-7}$。最细层记录进一步明确：
PyQCU 的 $14$ 与 QUDA 的 $37$ 都是最外层 Krylov 迭代数，不是所有层迭代之和；
三层粗化把最粗层体积减少 $16$ 倍，是本次主要加速来源。QUDA 1.1.0 的上游
recursive MultiGrid 未启用 double precision，因此 c128 只报告 PyQCU 侧结果，
不伪造跨库加速比。多卡部分只报告辅助吞吐与 MPI 能力门禁，不把 legacy 独立问题
并发误称为分布式 Strict MultiGrid。
\end{abstract}

\tableofcontents

\section{任务、验收标准与结论}

本轮任务的验收标准为：

\begin{enumerate}
  \item 两侧 MultiGrid 均具备默认关闭、可显式开启的分层计时与残差记录；
  \item 记录各层关键阶段，并在可获得处记录迭代前后的 $r$ norm；
  \item 使用同一输入、同一 odd parity、同一停止判据完成公平对照；
  \item 至少覆盖三层，并扩展格子与精度；
  \item 若 PyQCU 加速不明显，则分析 QUDA 优势与 PyQCU 短板并持续优化；
  \item 最后给出可追溯到原始 JSON/trace 的报告并完成回归。
\end{enumerate}

核心结论如下。

\begin{table}[htbp]
\centering
\caption{本任务的最终判断与证据等级}
\small
\begin{tabularx}{\textwidth}{@{}lL{0.27\textwidth}YL{0.23\textwidth}@{}}
\toprule
结论 & 数值证据 & 等级 & 来源 \\
\midrule
两侧 trace 均可开关 &
关闭时无 trace 文件、无同步；开启时事件完整 &
确证 &
\pathref{apply_multigrid_strict.cu}, \pathref{mg_profile.h} \\
三层显著快于 QUDA &
$0.653995$ s vs. $2.865627$ s，$4.3817\times$ &
确证 &
\pathref{data/strict_vs_quda_formal_3l_20260916_final.json} \\
最细层迭代口径已澄清 &
PyQCU 为 FGMRES 外层 $14$；QUDA 为 GCR 外层 $37$ &
确证 &
两侧 trace 与源码 \\
三层加速来自粗层体积下降 &
两层最粗体积 $49152$，三层为 $3072$，下降 $16$ 倍 &
推导 &
几何定义 \\
c128 QUDA 对照 &
QUDA 1.1.0 未启用 double MG，故只报告 PyQCU 侧 &
未验证（跨库） &
\pathref{data/mg_matrix_20260916/small-c128-l3/benchmark.json} \\
\bottomrule
\end{tabularx}
\end{table}

\section{可开关的 trace 设计}

\subsection{PyQCU Strict CUDA trace}

PyQCU 的 trace 由 \code{PYQCU\_STRICT\_TRACE\_FILE} 控制。变量未设置时，
\code{StrictFgmresTrace} 不打开文件，也不分配额外 trace workspace；正常生产路径
不分叉、不增加同步。实现位置为
\pathref{cpp/cuda/qcu/src/apply_multigrid_strict.cu:28}。

当前 trace version 为 2，记录：

\begin{itemize}
  \item \code{solve_begin/initial_residual/iteration/restart_residual/solve_end}：
        外层 right-preconditioned FGMRES 的迭代与 restart 真残差；
  \item \texttt{stage outer level phase seconds}：prepare、pre-smoother、restrict、
        recursive solve、prolongation、residual recompute、post-smoother 与
        reconstruct 的耗时；
  \item \texttt{residual outer level phase rnorm relative}：只在 trace 开启时用
        fine MATPC 或 coarse MATPC 计算真实 $r=b-Mx$。
\end{itemize}

主要阶段位点为
\pathref{cpp/cuda/qcu/src/apply_multigrid_strict.cu:2454}、
\pathref{cpp/cuda/qcu/src/apply_multigrid_strict.cu:3033}、
\pathref{cpp/cuda/qcu/src/apply_multigrid_strict.cu:3050}、
\pathref{cpp/cuda/qcu/src/apply_multigrid_strict.cu:3078} 与
\pathref{cpp/cuda/qcu/src/apply_multigrid_strict.cu:3083}。

\subsection{QUDA trace}

QUDA 私有头 \pathref{refer/git-rep/quda/lib/mg_profile.h} 使用
\code{QUDA\_MG\_TRACE\_FILE} 作为开关。未设置时 \code{enabled()} 为 false，
不会打开文件、同步设备或执行残差规约。

\code{MG::operator()} 的时间事件包括
\pathref{refer/git-rep/quda/lib/multigrid.cpp:1204}
至 \pathref{refer/git-rep/quda/lib/multigrid.cpp:1305}：

\begin{center}
\small
\begin{tabularx}{0.97\textwidth}{@{}lY@{}}
\toprule
字段 & 含义 \\
\midrule
\code{prepare\_outer/inner/coarsest} & parity/full 场的 prepare \\
\code{pre\_smoother/post\_smoother} & V-cycle 前后平滑器 \\
\code{restrict/prolongate} & $R$ 与 $P$ \\
\code{coarse\_solver} & 递归粗层或最粗 Krylov 求解 \\
\code{coarse\_solver\_iterations} & 粗层 solver 的 iteration 增量 \\
\code{after\_pre\_smoother} & 可取得时的真实 residual norm \\
\code{after\_prolongation/cycle\_exit} & 可取得时的粗层修正后残差 \\
\bottomrule
\end{tabularx}
\end{center}

外层 GCR 在 \pathref{refer/git-rep/quda/lib/inv_gcr_quda.cpp:355} 与
\pathref{refer/git-rep/quda/lib/inv_gcr_quda.cpp:370} 记录 iterated residual
及 restart 真残差。\code{OuterScope} 保证嵌套粗层 GCR 不会覆盖最外层
iteration 编号。

\subsection{残差语义}

若当前层算符为 $M$、右端为 $b$、当前解为 $x$，所有新增残差事件都定义

\begin{equation}
  r=b-Mx,\qquad r_{\rm abs}=\norm{r}_2,\qquad
  r_{\rm rel}=\frac{\norm{r}_2}{\norm{b}_2}.
\label{eq:residual}
\end{equation}

PyQCU 的 fine 层使用 odd-Schur MATPC 算符；QUDA 使用
\code{QUDA\_MATPC\_ODD\_ODD}。两侧均设置 \code{target\_parity=1}，
因此最细层迭代数差异不是 even/odd 记录反转造成的。

\section{最细层迭代口径与 parity}

旧报告中的 PyQCU 最细层“总迭代数”只有 $11$，容易与多层累计工作量混淆。
本轮明确区分三个量：

\begin{table}[htbp]
\centering
\caption{迭代计数层级}
\small
\begin{tabularx}{0.96\textwidth}{@{}lL{0.28\textwidth}L{0.42\textwidth}@{}}
\toprule
计数 & 当前定义 & 记录位置 \\
\midrule
外层迭代 &
PyQCU：right-preconditioned FGMRES column；QUDA：GCR iteration &
\code{solve_end} / \code{outer_iteration} \\
层内固定扫描 &
pre/post smoother 的 $\nu_{\rm pre}$、$\nu_{\rm post}$ &
PyQCU stage 数与 QUDA \code{pre/post_smoother_iterations} \\
粗层 Krylov 迭代 &
每一层递归求解器自身 iteration 增量 &
两侧 coarse solver 事件 \\
\bottomrule
\end{tabularx}
\end{table}

在 $16\times32\times32\times48$、c64、三层、五次正式样本中，PyQCU
外层 FGMRES 为 $14$ 次，QUDA 外层 GCR 为 $37$ 次。$14$ 不是“所有层总和”，
而是一次 odd-Schur solve 的最外层 Krylov 次数。三层 V-cycle 中每次外层迭代
包含：

\begin{equation}
  x_{k+1}^{\rm out}
  = \mathrm{postMR}\left[
      x_k+\mathrm{prolong}\left(
        \mathrm{coarseVcycle}\left(
          \mathrm{restrict}\left(\mathrm{preMR}(r_k)\right)
        \right)
      \right)
    \right].
\label{eq:preconditioner}
\end{equation}

因此前置条件器强度、粗层求解质量与 stop criterion 都会改变最外层次数，而不会
改变“一次 preconditioner call 对应一次外层迭代”的事实。

\section{公平对照协议}

\begin{table}[htbp]
\centering
\caption{正式与辅助测试条件}
\small
\begin{tabularx}{0.98\textwidth}{@{}lL{0.20\textwidth}L{0.20\textwidth}Y@{}}
\toprule
配置 & 格子 & 精度/层级 & 说明 \\
\midrule
历史正式基线 & $16\cdot32\cdot32\cdot48$ & c64 / 2 &
已有同输入正式对照，用于比较优化前后 \\
正式三层 & $16\cdot32\cdot32\cdot48$ & c64 / 3 &
五次 steady、两 warmup、cold x0 \\
小型矩阵 & $8\cdot8\cdot8\cdot16$ & c64 / 2,3 &
快速迭代与 stage trace \\
中型矩阵 & $16\cdot16\cdot16\cdot16$ & c64 / 3 &
体积扩展验证 \\
c128 可达性 & $8\cdot8\cdot8\cdot16$ & c128 / 3 &
PyQCU 通过；QUDA 1.1.0 fail-closed \\
多卡辅助 & $16\cdot32\cdot32\cdot48$ & c64 / legacy 2 &
独立问题吞吐，不是 distributed strict solve \\
\bottomrule
\end{tabularx}
\end{table}

两侧使用相同的：

\begin{itemize}
  \item gauge、source/RHS 和 canonical full null vectors；
  \item $n_v=12$、coarse spin $2$、$E=24$、block $(2,2,2,2)$；
  \item c64 时 $\kappa=1/(2m+8)$、$m=0.05$、$c_{\rm sw}=1$；
  \item odd--odd MATPC、周期边界与 scaled true-residual 定义为
        $\norm{b-Dx}/\norm{b}$；
  \item cold $x_0$、无日志解时间与 warmup 外 steady median。
\end{itemize}

\section{性能结果}

\subsection{正式三层结果}

\begin{table}[htbp]
\centering
\caption{正式大格 c64：优化前后与三层结果}
\small
\begin{tabular}{@{}lrrrrr@{}}
\toprule
配置 & PyQCU 中位/s & QUDA 中位/s & QUDA/PyQCU & 外层迭代 & 真残差 \\
\midrule
历史两层 & $1.966795$ & $2.094764$ & $1.0651$ & $11/37$ &
$3.60\cdot10^{-7}/7.30\cdot10^{-7}$ \\
当前三层 & $0.653995$ & $2.865627$ & $\mathbf{4.3817}$ & $14/37$ &
$4.90135\cdot10^{-7}/7.36873\cdot10^{-7}$ \\
\bottomrule
\end{tabular}
\end{table}

同一正式运行的 cache-hit setup 时间为 PyQCU $15.322$ s、QUDA $213.608$ s；
后者主要包含 QUDA 自身的粗算子构造与首次运行前准备。加速比只使用
steady solve-only median，setup 时间单列为工程可用性证据。

\begin{figure}[htbp]
\centering
\begin{tikzpicture}
\begin{axis}[
  width=0.82\textwidth,height=0.38\textwidth,
  xlabel={样本},ylabel={steady solve time / s},
  xtick={1,2,3,4,5},ymin=0,
  legend style={at={(0.5,1.03)},anchor=south,legend columns=2},
  grid=major]
\addplot+[mark=o,thick,color=pyblue] coordinates {
(1,0.6558688910008641)(2,0.6539949979924131)(3,0.653995474007111)(4,0.6566627829888603)(5,0.6509218419960234)};
\addplot+[mark=square*,thick,color=qudared] coordinates {
(1,2.8922347289917525)(2,2.856720073003089)(3,2.8656266540056095)(4,2.831856988996151)(5,3.019667030996061)};
\legend{PyQCU 三层,QUDA 三层}
\end{axis}
\end{tikzpicture}
\caption{正式大格 c64 五次 steady solve；数据来自项目数据目录中的正式 JSON。}
\end{figure}

三次扩展矩阵显示同一方向：

\begin{table}[htbp]
\centering
\caption{小型与中型 c64 三层 smoke 矩阵}
\small
\begin{tabular}{@{}lrrrrr@{}}
\toprule
配置 & PyQCU/s & QUDA/s & QUDA/PyQCU & 外层迭代 & 真残差上限 \\
\midrule
$8^3\cdot16$, 2L & $0.054420$ & $0.306282$ & $5.628$ & $9/32$ &
$7.97\cdot10^{-7}$ \\
$8^3\cdot16$, 3L & $0.045367$ & $0.951094$ & $20.96$ & $10/32$ &
$7.97\cdot10^{-7}$ \\
$16^3\cdot16$, 3L & $0.146060$ & $0.970265$ & $6.643$ & $31/58$ &
$6.89\cdot10^{-7}$ \\
\bottomrule
\end{tabular}
\end{table}

\subsection{分层 trace}

\begin{table}[htbp]
\centering
\caption{小型三层 trace 中的代表性阶段中位耗时；包含嵌套的位点明确标为 inclusive}
\small
\begin{tabular}{@{}llrr@{}}
\toprule
侧 & 阶段 & 中位/ms & 语义 \\
\midrule
PyQCU & level 0 \code{outer\_iteration} & $12.286$ & inclusive outer Krylov iteration \\
PyQCU & level 1 \code{coarse\_vcycle} & $8.098$ & inclusive recursive V-cycle \\
PyQCU & level 2 \code{coarsest\_bicgstab} & $2.113$ & coarsest solve \\
PyQCU & level 0 fine pre-smoother & $0.185$ & leaf stage \\
PyQCU & level 1 pre-smoother & $0.222$ & leaf stage \\
PyQCU & level 1 post-smoother & $0.198$ & leaf stage \\
QUDA & level 0 \code{coarse\_solver} & $30.862$ & inclusive recursive preconditioner \\
QUDA & level 1 \code{coarse\_solver} & $4.589$ & inclusive coarse solve \\
QUDA & level 0 pre-smoother & $0.533$ & leaf stage \\
QUDA & level 1 pre-smoother & $0.525$ & leaf stage \\
QUDA & level 0 restrict & $0.0912$ & leaf stage \\
QUDA & level 0 prolongate & $0.0829$ & leaf stage \\
\bottomrule
\end{tabular}
\end{table}

\subsection{$r$ 残差覆盖}

在小型三层 trace 中，PyQCU 每次外层迭代记录 fine/coarse 层的真实残差；
QUDA 记录 level 0/1 的 cycle enter、pre-smoother 后与可获得的后处理位置。
例如一个典型 PyQCU cycle：

\begin{center}
\small
\begin{tabular}{@{}llrr@{}}
\toprule
层 & 位点 & $\norm{r}_2$ & 相对值 \\
\midrule
0 & fine pre-smoother 后 & $0.29835$ & $0.29835$ \\
1 & level pre-smoother 后 & $0.09191$ & $0.36843$ \\
2 & coarsest BiCGStab 后 & $6.0796\cdot10^{-5}$ & $6.3087\cdot10^{-4}$ \\
1 & post-smoother 后 & $0.029622$ & $0.11874$ \\
\bottomrule
\end{tabular}
\end{center}

对应 QUDA 的代表事件为：

\begin{center}
\small
\begin{tabular}{@{}llrr@{}}
\toprule
层 & 位点 & $\norm{r}_2$ & 相对值 \\
\midrule
0 & cycle enter & $61.0568$ & $1$ \\
0 & pre-smoother 后 & $17.8278$ & $0.29199$ \\
1 & cycle enter & $14.5785$ & $1$ \\
1 & pre-smoother 后 & $5.44666$ & $0.37361$ \\
\bottomrule
\end{tabular}
\end{center}

\section{优化与原因分析}

\subsection{主要瓶颈}

历史两层 trace 中，PyQCU 每次外层迭代约 $179$ ms，其中最粗层 BiCGStab
约 $140$--$151$ ms；即约 $78\%$ 的最外层成本用于一个粗体积
$8\cdot16\cdot16\cdot24=49152$ 的序列求解。QUDA 同一阶段虽然也包含递归
粗层求解，但其 level 0 \code{coarse\_solver} 中位约 $115.8$ ms，且外层
iteration 更慢，因此总体与 PyQCU 两层同量级，存在细节差异但不是 parity
记录错误。

\subsection{三层为何有效}

加入第三层后，最粗层为

\begin{equation}
  4\cdot8\cdot8\cdot12=3072,
  \qquad
  \frac{49152}{3072}=16.
\label{eq:coarse-volume}
\end{equation}

正式大格的三层 PyQCU trace 中，最粗 BiCGStab 的中位耗时约
$4.89$ ms，而两层时约 $140$--$151$ ms。外层总成本从约 $179$ ms 降至
约 $78.8$ ms，外层迭代从 $11$ 增至 $14$，但总时间仍降至
$0.6545$ s。这符合多重网格的基本量纲：增加正确粗层会降低最粗 Krylov 的
体积和条件数负担，增加少量 transfer 成本换取更大的单层求解收益。

\subsection{实现改动}

\begin{itemize}
  \item \pathref{pyqcu/solver/_quda_multigrid.py} 新增
        \code{propagate\_null\_vectors}：后续 transition 的 $V$ 由上一层
        $R$ 作用到同一组 canonical basis，匹配 QUDA
        \code{generate\_all\_levels=false} 的粗基传播策略；专门测试见
        \pathref{examples/pyqcu/test_quda_multigrid.py}。
  \item \pathref{examples/qcu/dev87/bench_strict_vs_quda.py} 扩展
        \code{--lattice/--levels/--block/--gauge-path/--nullvec-path}，
        并支持三层 runtime cache。
  \item QUDA worker 仅为 transition 0 注入 QIO；后续层设置
        \code{num\_setup\_iter=0} 与 \code{generate\_all\_levels=false}，
        防止 QUDA 暗自生成另一套粗向量。
  \item \pathref{examples/qcu/dev87/bench_mg_matrix.py} 提供可复现矩阵；
        \pathref{examples/qcu/dev87/summarize_mg_trace.py} 聚合分层阶段与残差。
\end{itemize}

\section{多卡、精度与边界}

\subsection{多卡}

Strict 路径当前是单 rank fail-closed：C++ capability 的 production 多 rank
求解仍未启用，\code{fused\_fgmres=False}。因此在“一个 MPI job 分多个 rank
求解同一个 Strict MG”这一语义下，不能报告多卡求解加速。

可行的辅助测试是 legacy \code{MultiGpuMultigrid} 的独立问题并发：

\begin{center}
\small
\begin{tabular}{@{}lrrr@{}}
\toprule
配置 & 中位 wall/s & 单卡/双卡比 & 双卡效率 \\
\midrule
单 P100 独立完整问题 & $1.548730$ & $0.9591$ & $0.4796$ \\
双 P100 独立完整问题 & $1.614749$ & & \\
\bottomrule
\end{tabular}
\end{center}

结果说明双卡吞吐没有随卡数提升；这不是分布式 MG 的域分解结果。多 rank
strict preflight 与全局归约测试通过（$14$ passed / $13$ skipped，单 rank
测试在 2 rank 下按设计跳过）。需要真正分布式 Strict MG 时，必须补充
fine/compact/full halo、R/P coarse halo 和 distributed fused FGMRES。

\subsection{c128}

PyQCU c128 三层 $8^3\cdot16$ 在调整到可达的 fp64 相对残差门限后收敛：

\begin{center}
\small
\begin{tabular}{@{}lrrr@{}}
\toprule
项 & PyQCU & QUDA & 含义 \\
\midrule
steady median & $0.079139$ s & 不可用 & QUDA 未编译 double MG \\
外层迭代 & $14$ & --- & FGMRES outer \\
相对真残差 & $1.74532\cdot10^{-8}$ & --- & 达到 $5\cdot10^{-8}$ gate \\
\bottomrule
\end{tabular}
\end{center}

QUDA 1.1.0 源码默认关闭 \code{GPU\_MULTIGRID\_DOUBLE}；尝试强制构建
double MG 时，MMA double 模板不完整，故最终恢复为可复现的单精度 MG 构建。
这是明确的未验证项，不是把 PyQCU 结果冒名为对照结果。

\section{复现命令与证据}

\begin{verbatim}
source ./env.sh
source examples/qcu/dev87/quda_env.sh

python -B examples/qcu/dev87/bench_strict_vs_quda.py \
  --profile formal --levels 3 --repeats 5 --cache-expect hit \
  --side both --lattice 16 32 32 48 --precision c64 \
  --gauge-path data/gauge_16x32x32x48_m0.05_seed42_c64.h5 \
  --nullvec-path data/L16x32x32x48_nvec12_full_c64.h5 \
  --quda-nullvec-prefix data/L16x32x32x48_nvec12_quda \
  --quda-nullvec-manifest data/L16x32x32x48_nvec12_quda.conversion.json \
  --output data/strict_vs_quda_formal_3l_20260916_final.json

QUDA_MG_TRACE_FILE=/tmp/quda_mg.tsv \
PYQCU_STRICT_TRACE_FILE=/tmp/pyqcu_mg.tsv \
python examples/qcu/dev87/bench_strict_vs_quda.py ...
\end{verbatim}

\begin{longtable}{@{}L{0.28\textwidth}L{0.66\textwidth}@{}}
\caption{报告证据文件}\label{tab:evidence}\\
\toprule 文件 & 内容\\\midrule
\pathref{data/strict_vs_quda_formal_3l_20260916_final.json} &
正式三层 PyQCU/QUDA 五次 steady、真残差、内存和 provenance。\\
\pathref{data/mg_small_trace_20260916_r2.json.gz} &
小型三层逐外层、逐层 stage 与 residual trace。\\
\pathref{data/mg_small_trace_summary_20260916_r2.json} &
trace 聚合稳定值与 CSV 来源。\\
\pathref{data/mg_matrix_20260916/matrix_smoke.json} &
小/中型 c64 二、三层快速矩阵。\\
\pathref{data/mg_matrix_20260916/small-c128-l3/benchmark.json} &
PyQCU c128 收敛与 QUDA 不可用限制。\\
\pathref{data/multigpu_20260916.json} &
单 P100/双 P100 独立问题吞吐与一致性。\\
\bottomrule
\end{longtable}

\section{结论与后续}

本轮已经完成可开关 trace、三层公平对照和显著加速。当前最值得继续做的是：

\begin{enumerate}
  \item 将 QUDA 上游 double MG 的 MMA 缺口单独处理后，再补 c128 跨库正式对照；
  \item 若目标是分布式 Strict MG，按 MPI 阶段 1 的硬门禁依次实现 halo 与
        distributed fused FGMRES，再报告多卡加速；
  \item 继续以小型 smoke 迭代，只有冻结证据时才运行大格正式五次样本。
\end{enumerate}

\end{document}
